\chapter{Experimento Preliminar}
\label{ch:Experimento}

En esta sección se especifica la estrategia planificada para realizar un 
experimento preliminar que permita validar la metodología propuesta para 
la emulación física y extracción de métricas de rendimiento en 
arquitecturas de lógica reconfigurable. Se escaló el experimento a un 
tamaño pequeño para asegurar su viabilidad en el tiempo del proyecto y 
su funcionamiento correcto, lo que permitirá realizar ajustes a la 
metodología antes de su aplicación a un sistema más complejo. El 
principal interés del experimento es garantizar la correcta síntesis, 
mapeo, reteo y generación del bitstream de la arquitectura lógica 
reconfigurable mediante un circuito digital de prueba básico sin 
intervención de hardware físico.

\section{Metodología del experimento}

Se diseñará mediante código RTL una microarquitectura de lógica 
programable por un solo grupo de cuatro tablas de búsqueda, 
matrices de conmutación y su ruteo respectivo, todo con sus registros 
para asegurar que sea reconfigurable.

Como circuito de referencia para emular dentro de la microarquitectura 
reconfigurable descrita, se utilizará un multiplexor de dos entradas, 
se describirá de manera estructural compuerta por compuerta como se 
muestra en la función lógica \ref{eq1}. Dicho multiplexor es compuesto 
por dos compuertas AND, una compuerta NOT y una compuerta OR, cada una 
de las compuertas será asignada a una de las LUTs disponibles en la 
microarquitectura reconfigurable.

\begin{equation}
	\label{eq1}
	Y = (A \cdot \overline{S}) + (B \cdot S)
\end{equation}

Fidnalmente se realizará una simulación funcional y de temporización 
utilizando un banco de pruebas en software para verificar que las 
señales de entrada conmuten correctamente hacia la salida según el 
estado de la línea de selección.

\section{Resultados Esperados}

Por este medio se pretende:

\begin{itemize}

	\item Confirmar que las herramientas abiertas son capaces de 
	reconocer la estructura de la microarquitectura reconfigurable y 
	mapear la lógica combinacional del multiplexor.

	\item Verificar que el software puede interconectar las LUTs de 
	manera exitosa para propagar la señal del multiplexor.

	\item Obtener los primeros reportes de retardo estimados por el 
	software, sirviendo como línea base para la metodología de 
	caracterización final.

\end{itemize}
